\subsection{Sharp deletion order after averaging the pin slats}
\label{selection:sharp-section}

The preceding half-potential estimate is uniform at every pin. Averaging
over the actual deterministic vertical pin slats allows a smaller deletion
below \(u=4/3\). The result here concerns the same finite family, including
arbitrary probabilities inside its microscopic cells. No random shifts
or equidistribution of phases are assumed.

\begin{theorem}[Sharp averaged finite-model selection]
\label{selection:sharp-theorem}
Assume the finite model of Section~\ref{selection:finite-model}, with
\(0<\alpha<1/3\), and mollify at any fixed resolution
\(\epsilon\asymp\delta\). For \(Ca\le\rho\le r_0\), retain,
at a pin in track \(j\), exactly the source tracks satisfying
\( |c_i-c_j|\ge\rho\). Let \(f_y\) be the resulting mollified
pinned density. The deleted source mass is at most \(C\rho^\alpha\),
and
\begin{equation}\label{selection:sharp-bound}
 \int\|f_y\|_2^2\,d\nu_\delta(y)
 \le C\bigl(1+mr_0\rho^{\alpha-1}\bigr).
\end{equation}
The constants are independent of \(\delta\) and \(\rho\).
Consequently the cutoff
\begin{equation}\label{selection:optimal-cutoff}
 \rho_{\rm opt}=(mr_0)^{1/(1-\alpha)}
       \asymp r_0^{2/(1-\alpha)}
\end{equation}
gives a uniformly bounded averaged squared \(L^2\) norm with deleted
mass at most
\begin{equation}\label{selection:sharp-budget}
 C(mr_0)^{\alpha/(1-\alpha)}
       \asymp C\delta^{2\alpha\beta/(1-\alpha)}.
\end{equation}
This is the optimal order, up to constants depending on the norm bound,
for arbitrary positive selections with bounded averaged squared norm,
by the necessary estimate \eqref{selection:budget}.
\end{theorem}

\begin{proof}
Fix a pin track \(j\) and use the \emph{center-to-center} offsets
\(r_i=|c_i-c_j|\). Pins in its slat \(\ell\) are within \(O(w)\)
horizontally and \(O(\delta)\) vertically of \((c_j,3+\ell h)\).
For every retained track \(r_i\ge\rho\gg w\), put
\[
 t_{i,k,\ell}=\sqrt{r_i^2+(3+(\ell-k)h)^2}.
\]
The cell count of Lemma~\ref{selection:one-track} gives the uniform
slat envelope
\begin{equation}\label{selection:slat-envelope}
 f_{i,y}(t)\le
 C\frac{mh}{r_iw}\sum_k
 \mathbf1_{\{|t-t_{i,k,\ell}|\le Cr_iw\}}
 \quad(y\hbox{ in pin track }j,\hbox{ slat }\ell).
\end{equation}
To verify uniformity in the microscopic pin position, changing either
horizontal coordinate by \(O(w)\) changes distance by
\(O(r_iw+w^2)=O(r_iw)\). Vertical cell errors and mollification add
\(O(\delta)\), which is absorbed since \(\delta=w^2\).
The horizontal radial derivative has magnitude comparable to \(r_i\),
so at most \(C/r_i\) cells in a source slat contribute at one radial
value. Their masses are \(\asymp mh\delta/w\) and the mollifier's
supremum is \(O(\delta^{-1})\), proving the displayed amplitude.
This argument does not condition the vertical slat distribution on the
exact horizontal coordinate of a pin.

Tracks with \(r_i\ge r_0\) have total density bounded pointwise by a
constant, using Lemma~\ref{selection:one-track} and the comparison of
center offsets with actual offsets. Their total mass is at most one.
They contribute only a constant to the squared norm, including their
cross term with the remaining tracks. It suffices to handle
\(\rho\le r_i\le r_0\).

We first prove the following conditional averaged cross-track estimate.
For two such tracks, write
\(r=r_i\le R=r_{i'}\le r_0\) and \(\Delta=R^2-r^2\). Then
\begin{equation}\label{selection:cross-track}
 \frac1m\int_{y\ {
m in\ pin\ track}\ j}
       \int f_{i,y}(t)f_{i',y}(t)\,dt\,d\nu_\delta(y)
 \le Cm^2\left(1+\frac{r_0}{R}
               \mathbf1_{\{\Delta\le CRw\}}\right).
\end{equation}
The second term explicitly retains possible coherent equal-radius
resonances.

Multiply the envelopes \eqref{selection:slat-envelope}. An intersection
of their radial intervals has length at most \(Crw\), and is empty
unless the centers are within \(CRw\). Each pin slat has conditional
mass \(1/N_S\), so the left side of \eqref{selection:cross-track}
is at most
\begin{equation}\label{selection:lattice-count}
 C\frac{m^2h^2}{Rw}\frac1{N_S}
 \sum_\ell\sum_{k,k'}
 \mathbf1_{\{|t_{i,k,\ell}-t_{i',k',\ell}|\le CRw\}}.
\end{equation}
The use of a supremum within each pin slat justifies this estimate for
arbitrary microscopic cell probabilities.

Put \(q=k'-k\) and \(V=3+(\ell-k)h\). Both \(V\) and
\(V-qh\) lie in the fixed positive range of vertical source--pin
separations. Squaring the two central radii, their collision condition
implies exactly
\begin{equation}\label{selection:exact-phase}
 \big|\Delta-qh(2V-qh)\big|\le CRw.
\end{equation}
If \(q=0\), this requires \(\Delta\le CRw\). There are at most
\(N_S\) pairs \(k=k'\) per pin slat. Their contribution to
\eqref{selection:lattice-count} is at most
\[
 C\frac{m^2h^2}{Rw}N_S\,
       \mathbf1_{\{\Delta\le CRw\}}
 \le Cm^2\frac{r_0}{R}\mathbf1_{\{\Delta\le CRw\}}.
\]

Now fix \(q\ne0\) and \(k\). Increasing \(\ell\) by one
changes the expression inside the absolute value in
\eqref{selection:exact-phase} by exactly \(-2qh^2\). Thus at most
\[
 C\left(1+\frac{Rw}{|q|h^2}\right)
\]
pin slats can satisfy it. Any possible shift also obeys
\begin{equation}\label{selection:shift-range}
 |q|h\le C(\Delta+Rw)\le C(R^2+Rw),
\end{equation}
because \(2V-qh=V+(V-qh)\) is bounded below by a positive constant.
In particular
\[
 \frac{|q|h^2}{Rw}\le C(Rr_0+h)\le C.
\]
This absorbs the discrete counting error one. Since \(N_S\asymp h^{-1}\),
the fraction of pin slats satisfying the condition for each \(k,q\)
is at most \(CRw/(|q|h)\). Sum over \(k\), and substitute in
\eqref{selection:lattice-count}; this shift contributes at most
\(Cm^2/|q|\).

The possible nonzero shifts have bounded harmonic sum. If
\(\Delta>C_1Rw\), with a sufficiently large fixed \(C_1\),
\eqref{selection:exact-phase} forces \(q>0\) and
\[
 c\Delta\le qh\le C\Delta.
\]
These integers lie in an interval of fixed multiplicative ratio, so
\(\sum 1/q\lesssim1\), including when that interval meets one.
If \(\Delta\le C_1Rw\), then \eqref{selection:shift-range} and
\(Rw\le r_0w=h\) bound \(|q|\) by a fixed constant. The same
conclusion holds. Hence all nonzero shifts contribute at most \(Cm^2\),
proving \eqref{selection:cross-track}. This is a lattice count for the
actual pin slats, with no random-phase assumption.

It remains to sum the exceptional term. Its condition implies
\[
 |r_i-r_{i'}|\le Cw,
\]
since \(\Delta=(R-r)(R+r)\). For each fixed \(j,i\), only
\(O(1)\) centers \(c_{i'}\) satisfy this: they lie in intervals
of length \(O(w)\) around \(c_i\) or its reflection
\(2c_j-c_i\). The centers are \(3a\)-separated and \(w/a\to0\).
This counts exact reflected pairs as well as the diagonal \(i'=i\).
Their radii are comparable to \(r_i\), because \(r_i\gg w\).

The baseline terms in \eqref{selection:cross-track} sum to a constant,
since \(N_Tm=1\). The exceptional terms sum to at most
\[
 Cm^2r_0\sum_{r_i\ge\rho}r_i^{-1}.
\]
Dyadic annuli and \eqref{selection:center-count}, with \(\alpha<1\),
give
\[
 \sum_{r_i\ge\rho}m r_i^{-1}\le C_\alpha\rho^{\alpha-1}
 \qquad(\rho\ge Ca).
\]
Thus the conditional averaged squared norm in pin track \(j\) is at
most \(C(1+mr_0\rho^{\alpha-1})\), uniformly in \(j\). Averaging
these conditional estimates with track masses \(m\), and adding the
far-track contribution, proves \eqref{selection:sharp-bound}.

The removed mass is at most \(C(\rho^\alpha+a^\alpha)\), hence at
most \(C\rho^\alpha\). For \eqref{selection:optimal-cutoff}, use
\(m\asymp a^\alpha\asymp r_0\) to check
\[
 \rho_{\rm opt}/a\asymp
       a^{(3\alpha-1)/(1-\alpha)}\longrightarrow\infty,
 \qquad
 \rho_{\rm opt}/r_0\asymp
       r_0^{(1+\alpha)/(1-\alpha)}\longrightarrow0.
\]
The cutoff is therefore in the required range and makes
\(mr_0\rho^{\alpha-1}=1\). This proves
\eqref{selection:sharp-budget}; the necessary bound
\eqref{selection:budget} proves its asserted sharp order.
\end{proof}

The improvement in \eqref{selection:sharp-budget} concerns the joint
pin-averaged norm. Proposition~\ref{selection:finite-theorem} supplies
a larger deletion budget but a bound at \emph{every} pin; no improved
uniform pinwise bound is asserted here. The finite configuration varies
with \(\delta\). Consequently this theorem alone cannot be inserted
into Proposition~\ref{selection:compactness}, which requires restrictions
of one fixed pair of measures across the sequence of resolutions.
A persistent construction for a general fixed fractal measure remains a
separate missing step; no dimension-only distance theorem follows from
this finite-model result.
